\section{A countable completion with all principal parts}

The full principal parts also admit a Hilbert completion when
multiplicities are unbounded or zeros lie off the chosen line.
The construction below chooses the norm after the divisor and
finite interpolation functions have been specified.

Enumerate distinct zeros as $\rho_j$, $j\geq1$, with
$|\rho_j|\to\infty$, and let $r_j$ be their finite multiplicities.
Fix real $c$. Write $R_jf\in\mathbb C^{r_j}$ for the full residue
vector of $\mathcal Mf$ and set
\[
 z_j=i(c-\rho_j),\qquad K_j=z_jI-iS_j,
 \qquad (S_ja)_k=a_{k+1},\quad a_{r_j}=0.
\]
Each block has its Euclidean norm, so $\|S_j\|\leq1$.
For positive numbers $w_j$ define
\[
 \mathscr J_w=\left\{a=(a_j):
              \sum_jw_j\|a_j\|^2<\infty\right\},\qquad
 (K_wa)_j=K_ja_j.
\]
Its maximal domain is
\begin{equation}
 \operatorname{Dom}K_w=
 \left\{a\in\mathscr J_w:
             \sum_jw_j|z_j|^2\|a_j\|^2<\infty\right\}.
 \label{mo:full-domain}
\end{equation}
Indeed $K_w$ differs from multiplication by $z_j$ by the bounded
operator $-i\bigoplus_jS_j$. Equivalently its domain requires
$\sum_jw_j\|K_ja_j\|^2<\infty$.

\begin{theorem}[weighted full-jet completion]
\label{mo:all-jets}
For every such zero divisor and every real $c$, positive weights
$w_j$ can be chosen with the following properties. The map
$R:f\mapsto(R_jf)_j$ sends $\mathscr D$ into
$\operatorname{Dom}K_w$ and satisfies $RP_c=K_wR$.
Its range is dense and is a graph core for $K_w$.
The completion of $\mathscr D/\ker R$ in the norm $\|Rf\|$
is $\mathscr J_w$, and the induced dilation closes to $K_w$.
The graph of $R$ in $\mathscr H_c\oplus\mathscr J_w$ is dense.
Thus $R$ and its positive evaluation form are not closable in
the original norm.

The closed operator $K_w$ has compact resolvent and spectrum
$\{z_j:j\geq1\}$. The generalized eigenspace at $z_j$ is one
block of dimension $r_j$. It is self-adjoint exactly when all
$r_j=1$ and $\operatorname{Re}\rho_j=c$.
\end{theorem}
\begin{proof}
Let $I_n=\{(j,k):j\leq n,\ 0\leq k<r_j\}$ and denote its
coordinate basis by $e_h$. The finite graph theorem supplies,
for each $n$ and $h\in I_n$, a test $f_{n,h}$ such that
\begin{equation}
 \|f_{n,h}\|_{\mathscr H_c}\leq n^{-1},\qquad
 (R_jf_{n,h})_{j\leq n}=e_h.
 \label{mo:small-lifts}
\end{equation}
Here exact interpolation follows from graph density by correction
with the bounded finite-dimensional right inverse used in its proof.
All these choices precede the choice of weights.

For $g=U_cf$ supported in $[-m,m]$, each residue is an integral
against an exponential polynomial. Hence there is a finite
constant $C_{j,m}$ with
$\|R_jf\|\leq C_{j,m}\|g\|_2$ on that support.
One explicit choice is
\[
 C_{j,m}^2=\sum_{k=0}^{r_j-1}\int_{-m}^m
 \left|e^{(\rho_j-c)t}
 \sum_{\ell=0}^{r_j-k-1}
 b_{\rho_j,-k-\ell-1}\frac{t^\ell}{\ell!}\right|^2dt.
\]
Cauchy--Schwarz proves the bound. Define the finite maximum
\[
 A_j=\max\left(
 1,\ \max_{1\leq m\leq j}C_{j,m}^2,
 \max_{\substack{1\leq n<j\\h\in I_n}}
                  n^2\|R_jf_{n,h}\|^2\right),
 \qquad
 w_j=\frac{2^{-j}}{(1+|z_j|^2)A_j}.
\]
The last maximum is omitted for $j=1$.
For any fixed compact smooth test and sufficiently large $j$,
the summand $w_j(1+|z_j|^2)\|R_jf\|^2$ is at most
$2^{-j}\|U_cf\|_2^2$. This proves membership in the domain.
The coordinate identity from Proposition~\ref{mo:jet-comparison} therefore
holds as an identity in $\mathscr J_w$.

For a fixed coordinate $h$, take $n$ large enough that $h\in I_n$.
The first $n$ blocks of $Rf_{n,h}-e_h$ vanish. The remaining
blocks satisfy
\[
 \sum_{j>n}w_j(1+|z_j|^2)\|R_jf_{n,h}\|^2
 \leq n^{-2}\sum_{j>n}2^{-j}.
\]
Since $\|K_j\|\leq |z_j|+1$, this gives
$Rf_{n,h}\to e_h$ and $K_wRf_{n,h}\to K_we_h$.
Finite block truncations are a graph core for the maximal operator,
by convergence of the two sums defining its graph norm.
Thus $R\mathscr D$ is a graph core and is dense.
Equation \eqref{mo:small-lifts} also gives $f_{n,h}\to0$.
The graph closure contains $\{0\}\oplus\mathscr J_w$.
Subtracting its vertical vectors from $(f,Rf)$ and using density
of $\mathscr D$ proves density in the whole product.
The quotient and nonclosability assertions follow.

The maximal block operator is closed: convergence in its graph
implies convergence in every finite block, where each matrix is
closed. Its adjoint is the maximal operator with blocks
$\overline z_jI+iS_j^*$, as testing against finite block vectors
shows. These formulas also prove self-adjointness when all blocks
are real scalars. If it is symmetric, its restriction to every
finite block is symmetric. Theorem~\ref{mo:finite-symmetry} forces every
block to be a real scalar.

Choose $\lambda$ different from every $z_j$. For all sufficiently
large $j$, $|z_j-\lambda|>2$. The finite geometric series for
$(K_j-\lambda)^{-1}$ then gives
\[
 \|(K_j-\lambda)^{-1}\|
 \leq\frac1{|z_j-\lambda|-1}\longrightarrow0.
\]
The finitely many other inverses are bounded. Their direct sum
maps $\mathscr J_w$ into \eqref{mo:full-domain}, since
$K_ja_j=y_j+\lambda a_j$ for
$a_j=(K_j-\lambda)^{-1}y_j$.
Thus $\lambda$ belongs to the resolvent. Truncating this inverse
to finitely many finite-dimensional blocks converges in operator
norm, so the resolvent is compact. Each $z_j$ is an eigenvalue.
On every other block $K_l-z_j$ is invertible, so its generalized
eigenspace is precisely the $j$th block.
\end{proof}

The theorem requires no bound on the multiplicities. The norm
depends on the Laurent coefficients and the chosen interpolation
functions. Making a weight smaller preserves every estimate and
every graph approximation in the proof. For a divisor invariant
under $\rho\mapsto2c-\overline\rho$, one may therefore replace
the two weights on each orbit by their minimum. The resulting
residue involution is antiunitary. This does not make a nontrivial
Jordan block symmetric.

For a double zero on the line, the ordinary trace of a polynomial
in $K_j$ is $2p(\gamma_j)$, while the operator retains its
nonzero nilpotent part. More generally,
\[
 \operatorname{tr}p(K_j)=r_jp(z_j).
\]
This follows by triangularity. Consequently these traces cannot
distinguish the full block from $r_j$ repeated scalar eigenvalues.
Equality of spectral traces alone does not construct a map of
their full jet representations.
